\documentclass[ignorenonframetext,]{beamer}
\setbeamertemplate{caption}[numbered]
\setbeamertemplate{caption label separator}{: }
\setbeamercolor{caption name}{fg=normal text.fg}
\beamertemplatenavigationsymbolsempty
\usepackage{lmodern}
\usepackage{amssymb,amsmath}
\usepackage{ifxetex,ifluatex}
\usepackage{bm}
\usepackage{extarrows}
\usepackage{fixltx2e} % provides \textsubscript
\ifnum 0\ifxetex 1\fi\ifluatex 1\fi=0 % if pdftex
\usepackage[T1]{fontenc}
\usepackage[utf8]{inputenc}
\else % if luatex or xelatex
\ifxetex
\usepackage{mathspec}
\else
\usepackage{fontspec}
\fi
\defaultfontfeatures{Ligatures=TeX,Scale=MatchLowercase}
\fi
\usetheme{madrid}
\usecolortheme{orchid}
% use upquote if available, for straight quotes in verbatim environments
\IfFileExists{upquote.sty}{\usepackage{upquote}}{}
% use microtype if available
\IfFileExists{microtype.sty}{%
\usepackage{microtype}
\UseMicrotypeSet[protrusion]{basicmath} % disable protrusion for tt fonts
}{}
\newif\ifbibliography
\usepackage{natbib}
\usepackage{graphicx,grffile}
\makeatletter
\def\maxwidth{\ifdim\Gin@nat@width>\linewidth\linewidth\else\Gin@nat@width\fi}
\def\maxheight{\ifdim\Gin@nat@height>\textheight0.8\textheight\else\Gin@nat@height\fi}
\makeatother
% Scale images if necessary, so that they will not overflow the page
% margins by default, and it is still possible to overwrite the defaults
% using explicit options in \includegraphics[width, height, ...]{}
\setkeys{Gin}{width=\maxwidth,height=\maxheight,keepaspectratio}
\definecolor{azure}{rgb}{0, 0.4, 0.8}
\definecolor{crimson}{rgb}{0.8, 0.4, 0}
\definecolor{gold}{rgb}{0.72, 0.53, 0.04}
\definecolor{ciel}{rgb}{0.39, 0.58, 0.93}
\definecolor{dgreen}{rgb}{0.01, 0.75, 0.24}
% Prevent slide breaks in the middle of a paragraph:
\widowpenalties 1 10000
\raggedbottom

\setlength{\parindent}{0pt}
\setlength{\parskip}{6pt plus 2pt minus 1pt}
\setlength{\emergencystretch}{3em}  % prevent overfull lines
\providecommand{\tightlist}{%
\setlength{\itemsep}{0pt}\setlength{\parskip}{0pt}}
\setcounter{secnumdepth}{0}
\usepackage{bbm}
\usepackage{hhline}
\newcommand{\E}[1]{{\mathbb E}\left[#1\right]}
\newcommand{\Var}[1]{{\mathrm Var}\left(#1\right)}
\usepackage{beamerthemesplit}
\usepackage{amsmath}
\usepackage{amsthm}
\usepackage{kotex}
\usepackage{booktabs}
\usepackage{multirow}
\usepackage{makecell}
\usepackage{array}
\usepackage{arydshln}
\usepackage{caption}
\usepackage{comment}
\newcommand{\parallelsum}{\mathbin{\!/\mkern-5mu/\!}}
\newcommand\tab[1][1cm]{\hspace*{#1}}
\captionsetup[table]{belowskip=0pt,aboveskip=0pt}
\captionsetup[figure]{belowskip=0pt,aboveskip=0pt}
\newtheorem{model}{Model}

\title{Why do we need regression models for insurance?}

\author[Himchan Jeong]{\large{Himchan Jeong, PhD, FSA}}

%\institute{Simon Fraser University}

\date{}

\begin{document}
\frame{\titlepage}




\begin{frame}{Traditional Form of Insurance}

       \begin{columns}
          \column{0.5\linewidth}
             \centering
             \includegraphics[height=8cm, width=5cm]{shipwreck.jpg}
           \column{0.5\linewidth}
 \begin{itemize}
 \item \onslide<2-> If a seafarer was forced to throw cargo overboard to save the ship from sinking, the loss would be reimbursed collectively by his colleagues.
 \end{itemize}
   \end{columns} 
\end{frame}

\begin{frame}{Modern Form of Insurance}

       \begin{columns}
          \column{0.5\linewidth}
             \centering
             \includegraphics[height=8cm, width=5cm]{riskpooling.png}
           \column{0.5\linewidth}
 \begin{itemize}
 \item \onslide<2-> As modern society becomes much more complicated, risks need to be shared by larger pool of people, who do not necessarily know each other.
 \item \onslide<3-> As the risk pools have become anonymous, it has been critical to {\color{purple} assess and estimate} the amount of contribution to the pool.
 \end{itemize}
 \end{columns} 
\end{frame}

\begin{frame}{Example: Lung Cancer Insurance}
Suppose it is known that one out of 100 people is diagnosed with lung cancer every year. If 1000 people purchase lung cancer insurance that pays $\$10,000$ to whom is diagnosed with lung cancer, how much should the {\color{purple} annual premium per person} be?
 \begin{itemize}
 \item \onslide<2-> Expected total payment: $10,000 \times \frac{1,000}{100}=100,000$.
 \item \onslide<3-> Expected total payment per person: $\frac{100,000}{1,000}=100$.
 \end{itemize}
 \onslide<4-> \color{purple} Are we all set?
\end{frame}

\begin{frame}{Two Major Issues with Insurance Premium}
       \begin{columns}
          \column{0.4\linewidth}
 \begin{itemize}
 \item \onslide<2-> Is the expected total payment correct?
 \item \onslide<3-> Is it fair to charge the premium evenly?
 \end{itemize}
           \column{0.6\linewidth}
             \centering
 \onslide<4->    \includegraphics[height=6cm, width=7cm]{tobacco.jpg}
 \end{columns} 

\end{frame}

\begin{frame}{Example: Lung Cancer Insurance (Revisited)}
Suppose Company A has additional information that 
\begin{itemize}
	\item 1.5 out of 100 smokers, 
	\item 0.5 out of 100 non-smokers,
\end{itemize} 
are diagnosed with lung cancer every year. Further, 500 smokers and 500 non-smokers purchase lung cancer insurance that pays $\$10,000$ to whom is diagnosed with lung cancer.
 \begin{itemize}
 \item \onslide<2-> Expected total payment: $10,000 \times \frac{1.5 \times 500+0.5 \times 500}{100}=100,000$,
 \item \onslide<3-> Fair premium for a smoker: \onslide<4-> $10000 \times \frac{1.5}{100}=150$,
 \item \onslide<5-> Fair premium for a non-smoker: $10000 \times \frac{0.5}{100}=50$.
 \end{itemize}
\end{frame}

\begin{frame}{Example: Lung Cancer Insurance (Revisited)}
If Company A can differentiate the premium for a smoker and non-smoker,
\begin{table}[]
\begin{tabular}{l|rrrr}
\hline\hline
           & \# of Insured & Loss    & Premium & Difference \\ \hline
Smoker     & 500           & 75,000  & 75,000  & 0          \\
Non-smoker & 500           & 25,000  & 25,000  & 0          \\ \hline
Total      & 1,000         & 100,000 & 100,000 & 0        \\ \hline\hline
\end{tabular}
\end{table}

\end{frame}

\begin{frame}{Example: Lung Cancer Insurance (Revisited)}
If Company B cannot differentiate the premium for a smoker and non-smoker,
\begin{table}[]
\begin{tabular}{l|rrrr}
\hline\hline
           & \# of Insured & Loss    & Premium & Difference \\ \hline
Smoker     & 500           & 75,000  & 50,000  & -25,000          \\
Non-smoker & 500           & 25,000  & 50,000  & +25,000          \\ \hline
Total      & 1,000         & 100,000 & 100,000 & 0       \\  \hline\hline
\end{tabular}
\end{table}
\onslide<2-> After a year...
\begin{itemize}
	 \item \onslide<3-> 50 non-smokers left Company B to avoid higher premium (100 $\rightarrow$ 50),
 	 \item \onslide<4-> 50 smokers joined Company B to enjoy lower premium (150 $\rightarrow$ 100).
\end{itemize}

\end{frame}


\begin{frame}{Example: Lung Cancer Insurance (Revisited)}

Because of the change in the portfolio, Company B suffers $5,000$ of loss.
\begin{table}[]
\begin{tabular}{l|rrrr}
\hline\hline
           & \# of Insured & Loss    & Premium & Difference \\ \hline
Smoker     & 550           & 82,500  & 55,000  & -27,500          \\
Non-smoker & 450           & 22,500  & 45,000  & +22,500          \\ \hline
Total      & 1,000         & 105,000 & 100,000 & -5,000       \\  \hline\hline
\end{tabular}
\end{table}
\onslide<2-> And Company B decides to adjust the premium from 100 to 105 per person, to make up for the loss.
\end{frame}


\begin{frame}{Example: Lung Cancer Insurance (Revisited)}

After the adjustment in premium, the portfolio looks breakeven.
\begin{table}[]
\begin{tabular}{l|rrrr}
\hline\hline
           & \# of Insured & Loss    & Premium & Difference \\ \hline
Smoker     & 550           & 82,500  & 57,750  & -24,750          \\
Non-smoker & 450           & 22,500  & 47,250  & +24,750          \\ \hline
Total      & 1,000         & 105,000 & 105,000 & 0       \\  \hline\hline
\end{tabular}
\end{table}

\onslide<2-> However, after another year...
\begin{itemize}
	 \item \onslide<3-> 50 non-smokers left Company B to avoid higher premium (105 $\rightarrow$ 50),
 	 \item \onslide<4-> 50 smokers joined Company B to enjoy lower premium (150 $\rightarrow$ 105).
\end{itemize}
\end{frame}


\begin{frame}{Example: Lung Cancer Insurance (Revisited)}

Again, because of the change in the portfolio, Company B suffers $5,000$ of loss.
\begin{table}[]
\begin{tabular}{l|rrrr}
\hline\hline
           & \# of Insured & Loss    & Premium & Difference \\ \hline
Smoker     & 600           & 90,000  & 63,000  & -27,000          \\
Non-smoker & 400           & 20,000  & 42,000  & +22,000          \\ \hline
Total      & 1,000         & 110,000 & 105,000 & -5,000       \\  \hline\hline
\end{tabular}
\end{table}

\onslide<2-> \color{red} Would Company B survive in the end...?
\end{frame}


\begin{frame}{}

\begin{center}
\Huge \textbf{Insurance Data Analytics}
\end{center}

\end{frame}

\begin{frame}{Data Description}
Suppose we have a dataset of 1,000 observations described as follows:
$$
\begin{aligned}
\onslide<2->
Y:& \quad \text{Loss amount} \\
X_1:&\quad  \text{Gender} \ \ (1=\text{Male}, \ 0=\text{Female})\\
X_2:& \quad \text{Age}
\end{aligned}
$$

\onslide<3-> Our task is to predict the future loss amount given available information ({\color{blue} Gender} and {\color{blue} Age} in here).
\end{frame}


\begin{frame}{Remark: Machine Learning}
Note that machine learning can be classified into

\begin{itemize}
\item Supervised learning: task of learning a function that {\color{blue} maps an input to an output} based on example input-output pairs
\item  Unsupervised learning: type of algorithm that learns patterns from {\color{blue} untagged data}
\end{itemize}
In that sense, our task is one example of {\color{blue}`supervised learning'}.

\end{frame}


\begin{frame}{Supervised Learning}
             \centering
\includegraphics[height=9cm, width=12cm]{supervised.png}
\end{frame}

\begin{frame}{Unsupervised Learning}
             \centering
\includegraphics[height=9cm, width=12cm]{unsupervised.png}
\end{frame}



\begin{frame}{Model Selection}

\begin{itemize}
\item \onslide<1-> To predict the future loss, we need a model built upon observed data.
\item \onslide<1-> No model can capture reality exactly.
\item \onslide<1->Model is nothing but an approximation of reality with some assumptions.
\item \onslide<1->Some models help us to understand the complex dynamics of reality and predict the future outcomes.
\end{itemize}


\begin{quotation}
\onslide<2-> ``All models are wrong, but some models are useful.'' George Box
\end{quotation}
\end{frame}


\begin{frame}{Model vs Reality}

             \centering
\includegraphics[height=8cm, width=7cm]{modelreality.jpg}
\end{frame}

\begin{frame}{Preliminary Data Analysis}

\textbf{summary\color{blue}(Y)}

\begin{table}[]
\begin{tabular}{llllll}
Min. & 1st Qu. & Median & Mean & 3rd Qu. & Max. \\
2961 & 3337    & 3460   & 3454 & 3569    & 3901
\end{tabular}
\end{table}

\onslide<2-> Based on this data, we may propose our first model:
$$
\onslide<3-> \color{blue} \text{Model 1} \color{black}: P = 3454 \text{ for everyone}.
$$
\onslide<4-> {\color{red} Are you happy with it?}


\end{frame}


\begin{frame}{Preliminary Data Analysis}

 We might need a more refined model with risk classification.

       \begin{columns}
          \column{0.44\linewidth}
             \centering
   \onslide<2->    \includegraphics[height=5cm, width=8cm]{dimigo_scatter.pdf}
           \column{0.56\linewidth}
 	 \onslide<3-> One can see that the loss amounts tend to decrease as a function of {\color{blue} Age}.
 	 \end{columns}

\end{frame}


\begin{frame}{Effect of Age: Linear}

We may find the best line that captures linear effect of {\color{blue} Age} in relation to the loss amounts.

\begin{columns}
          \column{0.44\linewidth}
  \onslide<2->  \centering  \includegraphics[height=5cm, width=8cm]{dimigo_linear.pdf}
 	  \column{0.56\linewidth}
$$
\onslide<3-> \color{blue} \text{Model 2} \color{black}: P = 3618.204 - 3.815 X_2
$$
\onslide<4-> {\color{red} Are you happy with it?}
\end{columns}
\end{frame}


\begin{frame}{Effect of Age: Higher-order}

We may consider adding higher-order effect of {\color{blue} Age} in relation to the loss amounts.

\begin{columns}
          \column{0.44\linewidth}
  \onslide<2->  \centering  \includegraphics[height=5cm, width=8cm]{dimigo_higher.pdf}
 	  \column{0.56\linewidth}
$$
\begin{aligned}
\onslide<3-> & \color{blue}  \text{Model 3} \color{black}: \\
& P = 4154.728 - 31.838 X_2 + 0.324 X_2^2
\end{aligned}
$$
\onslide<4-> {\color{red} Looks better :)}

\onslide<5-> What about {\color{blue} Gender}?
\end{columns}
\end{frame}


\begin{frame}{More Preliminary Data Analysis}

We might plot the loss amounts for each gender with different colors.

       \begin{columns}
          \column{0.44\linewidth}
             \centering
   \onslide<2->    \includegraphics[height=5cm, width=8cm]{dimigo_gscatter.pdf}
           \column{0.56\linewidth}
 	 \onslide<3-> One can see that the level of loss amounts are also distinguishable with  {\color{blue} Gender}.
 	 \end{columns}

\end{frame}

\begin{frame}{Putting This Together}
\begin{columns}
          \column{0.44\linewidth}
  \onslide<1->  \centering  \includegraphics[height=5cm, width=8cm]{dimigo_final.pdf}
 	  \column{0.56\linewidth}
$$
\begin{aligned}
\onslide<2-> & \color{blue}  \text{Model 4} \color{black}: \\
& P = 3998.1628 - 30.3322 X_2 \\
& \qquad + 0.3074 X_2^2 + 201.8723X_1
\end{aligned}
$$

\end{columns}
\end{frame}


\begin{frame}{Validation}

\begin{table}[]
\begin{tabular}{l|rrrr}
\hline\hline              & Model 1 & Model 2 & Model 3 & Model 4 \\ \hline
Difference    & \onslide<2-> 0       & \onslide<3-> 0       & \onslide<4-> 0       & \onslide<4-> 0       \\
Average Error & \onslide<5-> 161.01  & \onslide<6-> 151.02  & \onslide<7-> 138.35  & \onslide<8-> 97.94   \\ \hline\hline
\end{tabular}
\end{table}

\begin{itemize}
	\item \onslide<9->   In terms of total loss prediction, all models are acceptable.
	\item \onslide<10-> By having a better risk classification model, however, one can predict the future loss amounts more accurately.
\end{itemize}

\end{frame}


\begin{frame}{Flood of Data}

\begin{columns}
          \column{0.5\linewidth}
    \centering  \includegraphics[height=5cm, width=5cm]{dataflood.jpg}
 	  \column{0.5\linewidth}
\begin{itemize}
	\item \onslide<2-> In this example, there were only two available inputs that can be used for risk classification, which allowed us to take (relatively simpler) step-wise approach.
	\item \onslide<3-> Due to advances in technology, nowadays we have much more data available compared to past years.
\end{itemize}
\end{columns}

\end{frame}


\begin{frame}{Issues Beyond Data Analytics}
\begin{itemize}
	\item  Regulations
		\begin{itemize}
			\item \onslide<1-> Use of gender for ratemaking has been banned in EU
		\end{itemize}
	\item \onslide<2-> Causality
				\begin{itemize}
			\item \onslide<3-> Shark attacks \textit{vs} Icecream sales
			\item \onslide<4-> Claims costs \textit{vs} Credit scores

		\end{itemize}

	\item \onslide<5-> Internal/external Communication
\end{itemize}
\end{frame}




\end{document}